% questions.tex —— 2012 年·广东卷·理科 · 题目内容片段（无答案）
% 由 tools/generate_projects.py 依来源总账自 2012/2012_数学真题_广东理_试题.md 生成，勿手改（改源 MD 后重跑）


\begin{gknotice}
\textbf{来源说明}：本题面依据固定源 PDF 的两页 A3 题面逐题转录，原卷印刷页码为 733—734。第 733 面为选择题 1—8、填空题 9—15，第 734 面为解答题 16—21（其中第 19 题排在该面左下）。

\textbf{图形说明}：第 6 题三视图、第 13 题程序框图、第 15 题圆与切线、第 17 题频率分布直方图、第 18 题四棱锥，均由项目内 TikZ 矢量源重绘；其余各题原页无图。

\textbf{解析归属}：固定源不附答案或解析；配套答案与解析均由本库按本题面独立复算。
\end{gknotice}

\gkpart{一、选择题}

\begin{question}
\gkti{1}设 $\mathrm i$ 为虚数单位，则复数 $\dfrac{5-6\mathrm i}{\mathrm i}=$\mbox{（　　）}

（A）$6+5\mathrm i$；（B）$6-5\mathrm i$；（C）$-6+5\mathrm i$；（D）$-6-5\mathrm i$。

\end{question}
\begin{question}
\gkti{2}设集合 $U=\{1,2,3,4,5,6\}$，$M=\{1,2,4\}$，则 $\complement_U M=$\mbox{（　　）}

（A）$U$；（B）$\{1,3,5\}$；（C）$\{3,5,6\}$；（D）$\{2,4,6\}$。

\end{question}
\begin{question}
\gkti{3}若向量 $\overrightarrow{BA}=(2,3)$，$\overrightarrow{CA}=(4,7)$，则 $\overrightarrow{BC}=$\mbox{（　　）}

（A）$(-2,-4)$；（B）$(3,4)$；（C）$(6,10)$；（D）$(-6,-10)$。

\end{question}
\begin{question}
\gkti{4}下列函数中，在区间 $(0,+\infty)$ 上为增函数的是\mbox{（　　）}

（A）$y=\ln(x+2)$；（B）$y=-\sqrt{x+1}$；（C）$y=\left(\dfrac12\right)^x$；（D）$y=x+\dfrac1x$。

\end{question}
\begin{question}
\gkti{5}已知变量 $x$，$y$ 满足约束条件
\nopagebreak[4]
\[\begin{cases}
y\leqslant 2,\\
x+y\geqslant 1,\\
x-y\leqslant 1,
\end{cases}\]
则 $z=3x+y$ 的最大值为\mbox{（　　）}

（A）$12$；（B）$11$；（C）$3$；（D）$-1$。

\Needspace{0.4\textheight}
\end{question}
\begin{question}
\gkti{6}某几何体的三视图如图所示，则它的体积为\mbox{（　　）}
\nopagebreak[4]
\begin{center}% ============================================================
% f2012_guangdong_science_q06_views.tex
% 2012 广东卷（理科）第 6 题 三视图（TikZ 矢量重绘）
% 正视图、侧视图：下方矩形 6×5，上方等腰三角形底 6、腰 5（高 4）
% 俯视图：直径 6 的圆；还原为圆柱（r=3, h=5）与圆锥（r=3, 母线 5, h=4）
% ============================================================
\begin{tikzpicture}[
  x=1cm, y=1cm, line width=0.6pt, font=\footnotesize, >=Stealth,
  dim/.style={<->, line width=0.4pt}
]

  % ---------------- 正视图 ----------------
  \begin{scope}
    \draw (-1.5,0) rectangle (1.5,2.5);
    \draw (-1.5,2.5) -- (0,4.5) -- (1.5,2.5);
    % 底边 6
    \draw[dim] (-1.5,-0.42) -- (1.5,-0.42);
    \draw (-1.5,-0.34) -- (-1.5,-0.50);
    \draw (1.5,-0.34) -- (1.5,-0.50);
    \node[below] at (0,-0.42) {$6$};
    % 高 5
    \draw[dim] (1.88,0) -- (1.88,2.5);
    \draw (1.80,0) -- (1.96,0);
    \draw (1.80,2.5) -- (1.96,2.5);
    \node[right] at (1.90,1.25) {$5$};
    % 两腰 5
    \node[rotate=53.13] at (-1.20,3.30) {$5$};
    \node[rotate=-53.13] at (1.20,3.30) {$5$};
    \node at (0,-1.10) {正视图};
  \end{scope}

  % ---------------- 侧视图 ----------------
  \begin{scope}[xshift=5.4cm]
    \draw (-1.5,0) rectangle (1.5,2.5);
    \draw (-1.5,2.5) -- (0,4.5) -- (1.5,2.5);
    \draw[dim] (-1.5,-0.42) -- (1.5,-0.42);
    \draw (-1.5,-0.34) -- (-1.5,-0.50);
    \draw (1.5,-0.34) -- (1.5,-0.50);
    \node[below] at (0,-0.42) {$6$};
    \draw[dim] (1.88,0) -- (1.88,2.5);
    \draw (1.80,0) -- (1.96,0);
    \draw (1.80,2.5) -- (1.96,2.5);
    \node[right] at (1.90,1.25) {$5$};
    \node[rotate=53.13] at (-1.20,3.30) {$5$};
    \node[rotate=-53.13] at (1.20,3.30) {$5$};
    \node at (0,-1.10) {侧视图};
  \end{scope}

  % ---------------- 俯视图 ----------------
  \begin{scope}[yshift=-4.5cm]
    \draw (0,0) circle (1.5);
    \fill (0,0) circle (1.1pt);
    \node at (0,-1.95) {俯视图};
  \end{scope}

\end{tikzpicture}
\end{center}
\nopagebreak[4]
（A）$12\pi$；（B）$45\pi$；（C）$57\pi$；（D）$81\pi$。

\end{question}
\begin{question}
\gkti{7}从个位数与十位数之和为奇数的两位数中任取一个，其个位数为 $0$ 的概率是\mbox{（　　）}

（A）$\dfrac49$；（B）$\dfrac13$；（C）$\dfrac29$；（D）$\dfrac19$。

\end{question}
\begin{question}
\gkti{8}对任意两个非零的平面向量 $\boldsymbol\alpha$ 和 $\boldsymbol\beta$，定义 $\boldsymbol\alpha\circ\boldsymbol\beta=\dfrac{\boldsymbol\alpha\cdot\boldsymbol\beta}{\boldsymbol\beta\cdot\boldsymbol\beta}$。若平面向量 $\boldsymbol a$，$\boldsymbol b$ 满足 $|\boldsymbol a|\geqslant|\boldsymbol b|>0$，$\boldsymbol a$ 与 $\boldsymbol b$ 的夹角 $\theta\in\left(0,\dfrac\pi4\right)$，且 $\boldsymbol a\circ\boldsymbol b$ 和 $\boldsymbol b\circ\boldsymbol a$ 都在集合 $\left\{\dfrac n2\mid n\in\mathbf Z\right\}$ 中，则 $\boldsymbol a\circ\boldsymbol b=$\mbox{（　　）}

（A）$\dfrac12$；（B）$1$；（C）$\dfrac32$；（D）$\dfrac52$。

\end{question}
\gkpart{二、填空题}

\begin{question}
\gkti{9}不等式 $|x+2|-|x|\leqslant 1$ 的解集为\gkblank{}。

\end{question}
\begin{question}
\gkti{10}$\left(x^2+\dfrac1x\right)^6$ 的展开式中 $x^3$ 的系数为\gkblank{}。（用数字作答）

\end{question}
\begin{question}
\gkti{11}已知递增的等差数列 $\{a_n\}$ 满足 $a_1=1$，$a_3=a_2^2-4$，则 $a_n=$\gkblank{}。

\end{question}
\begin{question}
\gkti{12}曲线 $y=x^3-x+3$ 在点 $(1,3)$ 处的切线方程为\gkblank{}。

\Needspace{0.4\textheight}
\end{question}
\begin{question}
\gkti{13}执行如图所示的程序框图，若输入 $n$ 的值为 $8$，则输出 $s$ 的值为\gkblank{}。
\nopagebreak[4]
\begin{center}% ============================================================
% f2012_guangdong_science_q13_flowchart.tex
% 2012 广东卷（理科）第 13 题 程序框图（TikZ 矢量重绘）
% 开始 → 输入 n → i=2, k=1, s=1 → 判断 i<n
%   是：s=(1/k)(s×i) → i=i+2 → k=k+1 → 回到判断
%   否：输出 s → 结束
% ============================================================
\begin{tikzpicture}[
  x=1cm, y=1cm, line width=0.6pt, font=\footnotesize, >=Stealth,
  every node/.style={inner sep=3pt},
  proc/.style={rectangle, draw, minimum width=2.9cm, minimum height=0.68cm},
  io/.style={trapezium, draw, trapezium left angle=70, trapezium right angle=110,
             minimum width=3.0cm, minimum height=0.72cm, trapezium stretches=true},
  dec/.style={diamond, draw, aspect=2.3, minimum width=2.5cm, minimum height=0.95cm},
  term/.style={rectangle, rounded corners=8pt, draw, minimum width=1.9cm, minimum height=0.60cm}
]

  \node[term] (start) at (0,0) {开始};
  \node[io]   (input) at (0,-1.25) {输入 $n$};
  \node[proc] (init)  at (0,-2.50) {$i=2$，$k=1$，$s=1$};
  \node[dec]  (judge) at (0,-3.95) {$i<n$};

  \node[proc, minimum height=1.10cm] (body) at (0,-5.55) {$s=\dfrac1k(s\times i)$};
  \node[proc] (inc)   at (0,-6.95) {$i=i+2$};
  \node[proc] (kk)    at (0,-8.20) {$k=k+1$};

  \node[io]   (out)   at (3.90,-5.55) {输出 $s$};
  \node[term] (end)   at (3.90,-6.95) {结束};

  \draw[->] (start) -- (input);
  \draw[->] (input) -- (init);
  \draw[->] (init) -- (judge);
  \draw[->] (judge) -- node[right] {是} (body);
  \draw[->] (body) -- (inc);
  \draw[->] (inc) -- (kk);
  \draw[->] (kk.west) -- (-2.55,-8.20) -- (-2.55,-3.95) -- (judge.west);
  \draw[->] (judge.east) -- (3.90,-3.95) -- (out.north);
  \node at (2.55,-3.72) {否};
  \draw[->] (out) -- (end);

\end{tikzpicture}
\end{center}
\nopagebreak[4]
\end{question}
\begin{question}
\gkti{14}在平面直角坐标系 $xOy$ 中，曲线 $C_1$ 和 $C_2$ 的参数方程分别为
\nopagebreak[4]
\[\begin{cases}
x=t,\\
y=\sqrt t,
\end{cases}\quad(t\text{ \text{为参数}})\]
\Needspace{5\baselineskip}
和
\nopagebreak[4]
\[\begin{cases}
x=\sqrt2\cos\theta,\\
y=\sqrt2\sin\theta,
\end{cases}\quad(\theta\text{ \text{为参数}}),\]
则曲线 $C_1$ 与 $C_2$ 的交点坐标为\gkblank{}。

\Needspace{0.4\textheight}
\end{question}
\begin{question}
\gkti{15}如图，圆 $O$ 的半径为 $1$，$A$，$B$，$C$ 是圆周上的三点，满足 $\angle ABC=30^\circ$，过点 $A$ 作圆 $O$ 的切线与 $OC$ 的延长线交于点 $P$，则 $PA=$\gkblank{}。
\nopagebreak[4]
\begin{center}% ============================================================
% f2012_guangdong_science_q15_circle.tex
% 2012 广东卷（理科）第 15 题 圆与切线（TikZ 矢量重绘）
% 圆 O 半径 1；A、B、C 在圆上，∠ABC=30°；
% 过 A 的切线与 OC 的延长线交于 P。图中 ∠AOC=60°，PA=√3。
% ============================================================
\begin{tikzpicture}[
  x=1cm, y=1cm, line width=0.6pt, font=\footnotesize, >=Stealth
]
  \coordinate (O) at (0,0);
  \coordinate (A) at (60:1.5);
  \coordinate (C) at (0:1.5);
  \coordinate (B) at (220:1.5);
  \coordinate (P) at (3.0,0);

  \draw (O) circle (1.5);
  \draw (A) -- (B) -- (C);
  \draw (O) -- (A);
  \draw (O) -- (P);
  \draw (A) -- (P);

  \fill (O) circle (1.0pt);
  \node[above left] at (O) {$O$};
  \node[above right] at (A) {$A$};
  \node[below left] at (B) {$B$};
  \node[below] at (C) {$C$};
  \node[right] at (P) {$P$};
\end{tikzpicture}
\end{center}
\nopagebreak[4]
\end{question}
\gkpart{三、解答题}

\begin{question}
\gkti{16}已知函数 $f(x)=2\cos\left(\omega x+\dfrac\pi6\right)$（其中 $\omega>0$，$x\in\mathbf R$）的最小正周期为 $10\pi$。

（1） 求 $\omega$ 的值；

（2） 设 $\alpha$，$\beta\in\left[0,\dfrac\pi2\right]$，$f\left(5\alpha+\dfrac{5\pi}{3}\right)=-\dfrac65$，$f\left(5\beta-\dfrac{5\pi}{6}\right)=\dfrac{16}{17}$，求 $\cos(\alpha+\beta)$ 的值。

\Needspace{0.4\textheight}
\end{question}
\begin{question}
\gkti{17}某班 $50$ 位学生期中考试数学成绩的频率分布直方图如图所示，其中成绩分组区间是：$[40,50)$，$[50,60)$，$[60,70)$，$[70,80)$，$[80,90)$，$[90,100]$。
\nopagebreak[4]
\begin{center}% ============================================================
% f2012_guangdong_science_q17_histogram.tex
% 2012 广东卷（理科）第 17 题 频率分布直方图（TikZ 矢量重绘）
% 组距 10；各组 [40,50)…[90,100] 的“频率/组距”依次为
% 0.006, 0.006, 0.010, 0.054, x, 0.006，其中 x=0.018
% ============================================================
\begin{tikzpicture}[
  x=1cm, y=1cm, line width=0.6pt, font=\footnotesize, >=Stealth
]
  \path[draw=none] (-0.75,-0.75) rectangle (8.05,7.15);
  % 坐标轴
  \draw[->] (-0.35,0) -- (7.55,0) node[below] {成绩};
  \draw[->] (0,-0.05) -- (0,6.55);
  \node[right, inner sep=3pt] at (0.10,6.05) {$\dfrac{\text{频率}}{\text{组距}}$};
  % 断轴符号
  \draw (-0.30,0) -- (-0.20,0.15) -- (-0.10,0) -- (0.00,0.15);

  % 六个矩形，高度按 100 倍纵坐标绘制
  \draw (0.5,0) rectangle (1.5,0.6);
  \draw (1.5,0) rectangle (2.5,0.6);
  \draw (2.5,0) rectangle (3.5,1.0);
  \draw (3.5,0) rectangle (4.5,5.4);
  \draw (4.5,0) rectangle (5.5,1.8);
  \draw (5.5,0) rectangle (6.5,0.6);

  % 虚线参考线
  \draw[dashed, line width=0.4pt] (0,5.4) -- (3.5,5.4);
  \draw[dashed, line width=0.4pt] (0,1.8) -- (4.5,1.8);
  \draw[dashed, line width=0.4pt] (0,1.0) -- (2.5,1.0);
  \draw[dashed, line width=0.4pt] (0,0.6) -- (5.5,0.6);

  % 纵轴刻度值
  \node[left] at (0,5.4) {$0.054$};
  \node[left] at (0,1.8) {$x$};
  \node[left] at (0,1.0) {$0.010$};
  \node[left] at (0,0.6) {$0.006$};

  % 横轴刻度
  \draw (0,0) -- (0,-0.09);
  \node[below] at (0,-0.10) {$0$};
  \foreach \xx/\lab in {0.5/40,1.5/50,2.5/60,3.5/70,4.5/80,5.5/90,6.5/100}{
    \draw (\xx,0) -- (\xx,-0.09);
    \node[below] at (\xx,-0.10) {$\lab$};
  }

\end{tikzpicture}
\end{center}
\nopagebreak[4]
（1） 求图中 $x$ 的值；

（2） 从成绩不低于 $80$ 分的学生中随机选取 $2$ 人，该 $2$ 人中成绩在 $90$ 分以上（含 $90$ 分）的人数记为 $\xi$，求 $\xi$ 的数学期望。

\Needspace{0.4\textheight}
\end{question}
\begin{question}
\gkti{18}如图所示，在四棱锥 $P-ABCD$ 中，底面 $ABCD$ 为矩形，$PA\perp$ 平面 $ABCD$，点 $E$ 在线段 $PC$ 上，$PC\perp$ 平面 $BDE$。

（1） 证明：$BD\perp$ 平面 $PAC$；

（2） 若 $PA=1$，$AD=2$，求二面角 $B-PC-A$ 的正切值。
\nopagebreak[4]
\begin{center}% ============================================================
% f2012_guangdong_science_q18_pyramid.tex
% 2012 广东卷（理科）第 18 题 四棱锥 P-ABCD（TikZ 矢量重绘）
% 底面 ABCD 为矩形，PA⊥平面 ABCD，E 在线段 PC 上，PC⊥平面 BDE
% 虚线为不可见棱：PA、AB、AD、AC、BD
% ============================================================
\begin{tikzpicture}[
  x=1cm, y=1cm, line width=0.6pt, font=\footnotesize, >=Stealth
]
  \coordinate (A) at (0,1.1);
  \coordinate (B) at (-2.5,0.1);
  \coordinate (C) at (2.1,0.1);
  \coordinate (D) at (4.6,1.1);
  \coordinate (P) at (-0.7,4.6);
  \coordinate (E) at ($(P)!0.5556!(C)$);

  % 不可见棱
  \draw[dashed, line width=0.45pt] (P) -- (A);
  \draw[dashed, line width=0.45pt] (A) -- (B);
  \draw[dashed, line width=0.45pt] (A) -- (D);
  \draw[dashed, line width=0.45pt] (A) -- (C);
  \draw[dashed, line width=0.45pt] (B) -- (D);

  % 可见棱
  \draw (P) -- (B);
  \draw (P) -- (C);
  \draw (P) -- (D);
  \draw (B) -- (C);
  \draw (C) -- (D);
  \draw (B) -- (E);
  \draw (E) -- (D);

  % 顶点与交点
  \fill (A) circle (1.0pt);
  \fill (B) circle (1.0pt);
  \fill (C) circle (1.0pt);
  \fill (D) circle (1.0pt);
  \fill (P) circle (1.0pt);
  \fill (E) circle (1.0pt);

  \node[above] at (P) {$P$};
  \node[below left] at (A) {$A$};
  \node[below left] at (B) {$B$};
  \node[below] at (C) {$C$};
  \node[right] at (D) {$D$};
  \node[above right] at (E) {$E$};

\end{tikzpicture}
\end{center}
\nopagebreak[4]
\end{question}
\begin{question}
\gkti{19}设数列 $\{a_n\}$ 的前 $n$ 项和为 $S_n$，满足 $2S_n=a_{n+1}-2^{n+1}+1$，$n\in\mathbf N^*$，且 $a_1$，$a_2+5$，$a_3$ 成等差数列。

（1） 求 $a_1$ 的值；

（2） 求数列 $\{a_n\}$ 的通项公式；

（3） 证明：对一切正整数 $n$，有 $\dfrac1{a_1}+\dfrac1{a_2}+\dfrac1{a_3}+\cdots+\dfrac1{a_n}<\dfrac32$。

\end{question}
\begin{question}
\gkti{20}在平面直角坐标系 $xOy$ 中，已知椭圆 $C:\dfrac{x^2}{a^2}+\dfrac{y^2}{b^2}=1$（$a>b>0$）的离心率 $e=\sqrt{\dfrac23}$，且椭圆 $C$ 上的点到点 $Q(0,2)$ 的距离的最大值为 $3$。

（1） 求椭圆 $C$ 的方程；

（2） 在椭圆 $C$ 上，是否存在点 $M(m,n)$，使得直线 $l:mx+ny=1$ 与圆 $O:x^2+y^2=1$ 相交于不同的两点 $A$，$B$，且 $\triangle OAB$ 的面积最大？若存在，求出点 $M$ 的坐标及相对应的 $\triangle OAB$ 的面积；若不存在，请说明理由。

\end{question}
\begin{question}
\gkti{21}设 $a<1$，集合 $A=\{x\in\mathbf R\mid x>0\}$，$B=\{x\in\mathbf R\mid 2x^2-3(1+a)x+6a>0\}$，$D=A\cap B$。

（1） 求集合 $D$（用区间表示）；

（2） 求函数 $f(x)=2x^3-3(1+a)x^2+6ax$ 在 $D$ 内的极值点。

\end{question}